\documentclass[10pt,a4paper]{amsart}
\usepackage[margin=19mm]{geometry}
\usepackage{amsmath,amssymb,mathtools}
\usepackage[scr=rsfs,cal=euler]{mathalfa}
\usepackage{xfrac}
\usepackage[unicode,hidelinks,bookmarks=false]{hyperref}
\newcommand{\ie}{i.e.\ }
\newcommand{\cf}{cf.\ }
\newcommand{\resp}{respectively\ }
\newcommand{\Subsection}{\subsection}
\providecommand{\nobreakdash}{\nobreak\hbox{-}}
\setlength{\emergencystretch}{2em}
\multlinegap0pt
\usepackage{mathtools}
\usepackage{bm}
\usepackage[shortlabels]{enumitem}
\usepackage{amssymb}
\usepackage{tikz}
\newtheorem{lemma}{Lemma}
\newtheorem{theorem}{Theorem}
\newtheorem{claim}{Claim}
\newcommand{\FORALL}[1]{\forall {#1} \, }
\newcommand{\IMPLIES}{\Rightarrow }
\newcommand{\Mid}{\boldsymbol\mid}			% separator
\newcommand{\On}{\mathord{\mathrm{Ord}}}		% the class of ordinals
\newcommand{\R}{\mathbb{R}}	% the reals
\newcommand{\card}[1]{\mathopen | #1 \mathclose |}		% cardinality of a set
\newcommand{\cardLR}[1]{ \left | #1 \right |}		% cardinality of a set
\newcommand{\set}[1]{\mathopen \{ {#1} \mathclose \}} % set 
\newcommand{\setLR}[1]{\left \{ {#1} \right \}} % set with large delimiters
\newcommand{\setof}[2]{\mathopen \{{#1}\Mid{#2} \mathclose\}} % set of...
\title{How many sprays cover the space?}
\author{Alessandro Andretta, Ivan Izmestiev}
\date{Source: arXiv:2406.04078v2 (CC BY 4.0)}
\begin{document}
\maketitle

\noindent\emph{Source and licence.} How many sprays cover the space?. Version arXiv:2406.04078v2 (CC BY 4.0), distributed by arXiv under the Creative Commons Attribution 4.0 International licence, http://creativecommons.org/licenses/by/4.0/, as recorded in the arXiv licence metadata for this exact version and shown on the abstract page https://arxiv.org/abs/2406.04078v2. Full licence notice: \texttt{licenses/arxiv-2406.04078-CC-BY-4.0.txt}.\par\smallskip

\noindent\textit{Editorial scope.} Counted unit: the article's theorem th:EJM-Rd (source0.tex lines 938--949). The article's complete original proof of it is delivered with it (source0.tex lines 1328--1397, one \texttt{proof} environment of 70 lines, which the article places away from the statement). No mathematics is written, repaired or re-derived: the statement and the proof are the article's own lines, in the article's own notation, and no retained line is substituted or re-flowed.  Retained uncounted as the source-local context the proof needs: lines 921--924; lines 1189--1199.  Nothing was omitted from any retained span, so the package carries no omission record.  No retained line contains a citation, so the wrapper carries no bibliography and no bibliography entry.  No retained line contains a figure environment, an image-inclusion command or drawing code, so no image is delivered and the package declares no asset dependency.  Editorial formatting only, in addition: an AMS \texttt{amsart} preamble that restates the article's own declarations and macros used by the retained lines, namely the environments \texttt{lemma}, \texttt{theorem}, \texttt{claim} on separate counters, together with the standard packages those lines need.  The retained environments print in this document as Lemma 1, Theorem 1, Claim 1 in order of appearance, whereas the article prints the same items with the numbering of its own full document.  The complete native source of the article as distributed in the arXiv e-print for this exact version is bundled unchanged as \texttt{sources/160399/source0.tex}.
\begin{lemma}\label{lem:setofdifferenceinabeliangroup}
Let \( \varepsilon > 0 \).
For all \( X \subseteq \R \) with \( \aleph_0 \leq \card{X} < \kappa \leq 2^{\aleph_0} \), there is \( S \subseteq ( - \varepsilon ; \varepsilon ) \) such that \( \card{S} = \kappa \) and \( ( X - X ) \cap ( S - S ) = \set{ 0 } \).
\end{lemma}

\begin{theorem}\label{th:hyperplanes=>bound}
Let \( \delta \in \On \), \( d \geq 3 \), \( n \geq 0 \), and let \( N = ( n + 1 ) ( d - 1 ) + 1 \).
Suppose \( \mathbf{u}_1 , \dots , \mathbf{u}_N \) are distinct, non-zero vectors of \( \R^d \), and that \( A_1 , \dots , A_ N \subseteq \R^d \) are such that for all \( \mathbf{p} \in \R^d \)
\[
 \FORALL{ 1 \leq k \leq N} \left ( \card{ H_{\mathbf{u}_k } ( \mathbf{p} ) \cap A_k } < \aleph_ \delta \right ) .
\]
If \( 2^{\aleph_0} > \aleph_{ \delta + n } \) then for every \( \varepsilon > 0 \) there is \( Z_{ n , \varepsilon } \subseteq ( - \varepsilon ; \varepsilon )^d \) of size \( \aleph_{ \delta + n + 1} \) such that 
\[ 
 \FORALL{ \mathbf{p} \in \R^d } \left ( \mathbf{p} + Z_{ n , \varepsilon } \nsubseteq \textstyle \bigcup_{ k = 1 }^N A_k \right ) . 
\]
\end{theorem}

\begin{theorem}\label{th:EJM-Rd}
Let \( \delta \in \On \), \( d \geq 3 \) and let \( N = 2 ( d - 1 ) \).
Suppose \( \mathbf{u}_1 , \dots , \mathbf{u}_N \) are distinct, non-zero vectors of \( \R^d \).
Suppose \( A_1 , \dots , A_N \) are subsets of \( \R^d \) such that for all \( \mathbf{p} \in \R^d \), and for all \( 1 \leq i \leq N \)
\[
\card{ H_{\mathbf{u}_i} ( \mathbf{p} ) \cap A_i } < \aleph_ \delta . 
\]
If \( 2^{\aleph_0 } > \aleph_{ \delta } \) then for every \( \varepsilon > 0 \) there is \( Z \subseteq ( - \varepsilon ; \varepsilon ) ^ d \) of size \( \aleph_{ \delta + 1 } \) such that 
\[
\FORALL{ \mathbf{p} \in \R^d } ( \mathbf{p} + Z \nsubseteq \textstyle \bigcup_{i = 1 }^N A_i ) . 
\]
\end{theorem}

\begin{proof}[Proof of Theorem~\ref{th:hyperplanes=>bound}]
Fix \( \delta \in \On \) and \( d \geq 3 \), and proceed by induction on  \( n \).

The case \( n = 0 \) is Theorem~\ref{th:EJM-Rd}, so assume the result holds for some \( \bar{n} \) towards proving it for \( \bar{n} + 1 \).

Suppose \( 2^{\aleph_0} > \aleph_{ \delta + \bar{n} + 1 } \).
Let \( N = ( \bar{n} + 2 ) ( d - 1 ) + 1 \), and let \( \mathbf{u}_1 , \dots , \mathbf{u}_ N \) and \( A_1 , \dots , A_ N \) be as in the statement.
Given \( \varepsilon > 0 \) we must construct \( Z = Z_{ \bar{n} + 1 , \varepsilon } \) such that 
\[ 
Z \subseteq ( - \varepsilon ; \varepsilon )^d , \quad \card{Z} = \aleph_{ \delta + \bar{n} + 2 } , \quad \FORALL{ \mathbf{p} \in \R^d } \bigl ( \mathbf{p} + Z \nsubseteq \textstyle \bigcup_{k = 1 }^N A_k \bigr ) .
\] 
Let \( \bar{N} = ( \bar{n} + 1 ) ( d - 1 ) + 1 \).
Then \( \mathbf{u}_1 , \dots , \mathbf{u}_{\bar{N} } \) and \( A_1 , \dots , A_{ \bar{N} } \) satisfy the hypotheses of the statement for \( \bar{n} \).
By inductive assumption there is \( \bar{Z} = Z_{ \bar{n} , \varepsilon / 2 } \subseteq ( - \varepsilon / 2 ; \varepsilon / 2 )^d \) of size \( \aleph_{ \delta + \bar{n} + 1 } \) such that 
\begin{equation}\label{eq:th:hyperplanes=>bound-1}
 \FORALL{ \mathbf{p} \in \R^d } \bigl ( \mathbf{p} + \bar{Z} \nsubseteq \textstyle \bigcup_{ k = 1 }^{\bar{N}} A_k \bigr ) .
\end{equation}
As \( N - \bar{N} = d - 1 \), there is a non-zero vector \( \mathbf{v} \) such that 
\begin{equation}\label{th:hyperplanes=>bound-2}
 \bar{N} < k \leq N \IMPLIES \mathbf{v} \boldsymbol{\cdot} \mathbf{u}_k = 0 . 
\end{equation}
The subspace \( V = \setof{ r \mathbf{v}}{ r \in \R } \) is of cardinality \( 2^{\aleph_0 } \geq \aleph_{ \delta + \bar{n} + 2 } \), and \( X = \setof{ r \in \R }{ r \mathbf{v} \in \bar{Z}} \) is of cardinality \( \aleph_{ \delta + \bar{n} + 1 } \), so by Lemma~\ref{lem:setofdifferenceinabeliangroup} we obtain \( S \subseteq ( - \varepsilon / 2 ; \varepsilon / 2 ) \) of size \( \aleph_{ \delta + \bar{n} + 2 } \) such that \( ( S - S ) \cap ( X - X ) = \set{0} \).
Letting \( Y = \setof{ s \mathbf{v} }{ s \in S } \) then \( \card{Y} = \aleph_{ \delta + \bar{n} + 2 } \) and \( Y \subseteq ( - \varepsilon / 2 ; \varepsilon / 2 )^d \) and 
\[
 ( Y - Y ) \cap ( \bar{Z} - \bar{Z} ) = \setLR{\mathbf{0}} .
\]
Let 
\[
 Z \coloneqq Y + \bar{Z}= \textstyle \bigcup_{ \mathbf{y} \in Y} \mathbf{y} + \bar{Z} = \bigcup_{ \mathbf{z} \in \bar{Z} } \mathbf{z} + Y .
\]
Observe that \( Z \subseteq ( - \varepsilon ; \varepsilon )^d \) and is of cardinality \( \aleph_{ \delta + \bar{n} + 2 } \).
We must argue that \( \mathbf{p} + Z \nsubseteq \bigcup_{ k = 1 }^ N A_k \) for all \( \mathbf{p} \in \R^d \).

Towards a contradiction, let \( \hat{\mathbf{p}} \in \R^d \) be such that \( \hat{\mathbf{p}} + Z \subseteq \bigcup_{ k = 1 }^N A_k \).

\begin{claim}\label{cl:th:hyperplanes=>bound-1}
If \( \mathbf{y} , \mathbf{y}' \in Y \) and \( ( \mathbf{y} + \bar{Z} ) \cap ( \mathbf{y}' + \bar{Z} ) \neq \emptyset \) then \( \mathbf{y} = \mathbf{y}' \).
\end{claim}

\begin{proof}
Suppose \( \mathbf{y} + \mathbf{z} = \mathbf{y}' + \mathbf{z}' \) with \( \mathbf{y} , \mathbf{y}' \in Y \) and \( \mathbf{z} , \mathbf{z}' \in \bar{Z} \).
Then \( \mathbf{y} - \mathbf{y} ' = \mathbf{z}' - \mathbf{z} \in ( Y - Y ) \cap ( \bar{Z} - \bar{Z} ) = \set{ \mathbf{0} } \), so \( \mathbf{y} = \mathbf{y}' \).
\end{proof}

Recall that \( Y \subseteq V \) and by~\eqref{th:hyperplanes=>bound-2} \( V \) is contained in \( H_{\mathbf{u}_k } ( \mathbf{0} ) \), for \( \bar{N} < k \leq N \).
By assumption \( \card{ A_k \cap H_{\mathbf{u}_k } ( \mathbf{q} ) } < \aleph_ { \delta } \), so \( \card{ A_k \cap ( \mathbf{q} + Y ) } < \aleph_ { \delta } \), for each \( \mathbf{q} \in \R^d \).
Therefore for any \( \bar{N} < k \leq N \),
\begin{equation}\label{th:hyperplanes=>bound-3}
\card{ A_k \cap ( \hat{\mathbf{p}} + Z ) } = \card{ A_k \cap ( \hat{\mathbf{p}} + Y + \bar{Z} ) } = \cardLR{ \textstyle \bigcup_{ \mathbf{z} \in \bar{Z} } A_k \cap ( \hat{\mathbf{p}} + \mathbf{z} + Y ) } \leq \aleph _{ \delta + \bar{n} + 1 } . 
\end{equation}

\begin{claim}\label{cl:th:hyperplanes=>bound-2}
There is \( \hat{\mathbf{y}} \in Y \) such that \( ( \hat{\mathbf{p}} + \hat{\mathbf{y}} + \bar{Z} ) \cap \bigcup_{ \bar{N} < k \leq N} A_k = \emptyset \).
\end{claim}

\begin{proof}
Towards a contradiction, suppose that for all \( \mathbf{y} \in Y \) there are \( \mathbf{z} ( \mathbf{y} ) \in \bar{Z} \) and \( \bar{N} < k ( \mathbf{y} ) \leq N \) such that \( \hat{\mathbf{p}} + \mathbf{y} + \mathbf{z} ( \mathbf{y} ) \in A_{ k ( \mathbf{y} ) } \).
As 
\[
 \aleph_{ \delta + \bar{n} + 2 } = \card{Y} > \aleph_{ \delta + \bar{n} + 1 } =\bar{Z} > \card{\setof{ k }{ \bar{N} < k \leq N}} 
\] 
there is \( \hat{Y} \subseteq Y \) of size \( \aleph_{ \delta + \bar{n} + 2 } \), and \( \hat{\mathbf{z}} \in \bar{Z} \) and \( \bar{N} < \hat{k} \leq N \) such that \( \mathbf{z} ( \mathbf{y} ) = \hat{\mathbf{z}} \) and \( k ( \mathbf{y} ) = \hat{k} \) for all \( \mathbf{y} \in \hat{Y} \).
Therefore \( \FORALL{ \mathbf{y} \in \hat{Y} } ( \hat{\mathbf{p}} + \mathbf{y} + \hat{\mathbf{z}} \in A_{\hat{k}} ) \).
By Claim~\ref{cl:th:hyperplanes=>bound-1} the map \( \hat{Y} \to A_{\hat{k}} \), \( \mathbf{y} \mapsto \hat{\mathbf{p}} + \mathbf{y} + \hat{\mathbf{z}} \) is injective, so \( \card{ A_{\hat{k}} \cap ( \hat{\mathbf{p}} + Y + \bar{Z} ) } \geq \card{ \hat{Y} } = \aleph _{ \delta + \bar{n} + 2 } \), against~\eqref{th:hyperplanes=>bound-3}.
\end{proof}

Fix \( \hat{\mathbf{y}} \) as in Claim~\ref{cl:th:hyperplanes=>bound-2}.
Then \( ( \hat{\mathbf{p}} + \hat{ \mathbf{y} } ) + \bar{Z} \subseteq A_1 \cup \dots \cup A_{\bar{N}} \) against~\eqref{eq:th:hyperplanes=>bound-1}.
Having reached a contradiction, we conclude that \( \mathbf{p} + Z \nsubseteq \bigcup_{ k = 1 }^ N A_k \) for all \( \mathbf{p} \in \R^d \).
\end{proof}

\end{document}
